\documentclass[10pt]{article}

\usepackage[utf8]{inputenc}

\usepackage[T1]{fontenc}

\usepackage{amsmath}

\usepackage{amsfonts}

\usepackage{amssymb}

\usepackage[version=4]{mhchem}

\usepackage{stmaryrd}

\usepackage{bbold}

\usepackage{graphicx}

\usepackage[export]{adjustbox}

\graphicspath{ {./images/} }



\begin{document}

согласно к-рой обе выборки извлечены из одной и той же совокупности. Если



\[

X_{(1)} \leqslant X_{(2)} \leqslant \ldots \leqslant X_{(n)} \text { и } Y_{(1)} \leqslant Y_{(2)} \leqslant \ldots \leqslant Y_{(m)}

\]



\begin{itemize}

  \item вариационные ряды, отвечающие данным выборкам, a $F_{n}(x)$ І $G_{m}(x)$ - соответствующие им функции эмпири'І. распределсния, то гипотөзу $H_{0}$ можно заиисать в виде следующего тождества

\end{itemize}



\[

H_{0}: \mathrm{EF}_{n}(x) \equiv \mathrm{E} G_{m}(x)

\]



Далее, пусть в качестве возможных альтернатив к $H_{0}$ рассматриваются гипотезы



\[

\begin{aligned}

& H_{1}^{+}=\sup _{|x|<\infty} \mathrm{E}\left[G_{m}(x)-F_{n}(x)\right]>0, \\

& H_{1}^{-}=\inf _{|x|<\infty} \mathrm{E}\left[G_{m}(x)-F_{n}(x)\right]<0, \\

& H_{1}=\sup _{|x|<\infty}\left|\mathrm{E}\left[G_{m}(x)-F_{n}(x)\right]\right|>0 .

\end{aligned}

\]



Для проверки гипотезы $H_{0}$ против односторонних альтернатив $H_{1}^{+}$и $H_{1}^{-}$, а также против двусторонней $H_{1}$, Н. В. Смирновым предложены статистич. критерии, ностроенные на статистиках



\[

\begin{gathered}

D_{m, n}^{+}=\sup _{|x|<\infty}\left[G_{m}(x)-F_{n}(x)\right]=\max _{1 \leqslant k \leqslant m}\left(\frac{k}{m}-F_{n}\left(X_{(k)}\right)\right), \\

D_{m, n}^{-}=-\inf _{|x|<\infty}\left[G_{m}(x)-F_{n}(x)\right]= \\

=\max _{1 \leqslant k \leqslant m}\left(F_{n}\left(X_{(k)}\right)-\frac{k-1}{m}\right), \\

D_{m, n}=\sup _{|x|<\infty}\left|G_{m}(x)-F_{n}(x)\right|=\max \left(D_{m, n}^{+}, D_{m, n}^{-}\right)

\end{gathered}

\]



соответственно, причем, как следует из определений статистик $D_{m, n}^{+}$и $D_{m}^{-}, n$, при справедливости проверяемой гипотезы $H_{0}$, статистики $D_{\bar{m}, n \text { и }}^{-} D_{m, n}^{+}$распределены одинаково. В основе этих критериев лежит следующая теорема: если $\min (m, n) \rightarrow \infty$ так, что отношение $m / n$ остается постоянным, то при справедливости гипотезы $H_{0}$ для любого $y>0$



\[

\begin{gathered}

\lim _{m \rightarrow \infty} \mathrm{P}\left\{\sqrt{\frac{m n}{m+n}} D_{m, n}^{+}<y\right\}=1-e^{-2 y^{2},}, \\

\lim _{n \rightarrow \infty} \mathrm{P}\left\{\sqrt{\frac{m n}{m+n}} D_{m, n}<y\right\}=K(y),

\end{gathered}

\]



где $K(y)$ - функция распределения Колмогорова. Были получены (см. [4] - [6]) асимптотич. разложения для функций распределения статистик $D_{m, n}^{+}$и $D_{m, n}$.



Согласно С. к. с уровнем значимости $Q$ гипотеза $H_{0}$ отвергается в пользу одной из рассматриваемых альтернатив $H_{1}^{+}, H_{1}^{-}, H$, если статистика, соответствующая выбранной альтернативе, превосходит $Q$ - критическое значение критерия, для вычисления к-рого рекомендуется пользоваться аппроксимациями, полученными J. Н. Большевым [2] с помощыо асимптотит. преобразований Пирсона.



Сі. также Колмогорова критерий, Колмогорова Смирнова критерий.



в. Лит:: [1] С м и р но в Н. В., "Бюлл. МГУ", 1939, т. 2, ее примен.», 1963 , т. 8, в. 2, с. $129-55$; [3] Б Б л ь ш е в Л. Н., С м ир но в H. В., Таблицы математической статистики, 2 изд., М., 1968; [4] К о р о л ю к В. С., «Теория вероятн. и ее примен.», 1959 , т. 4, в. 4, с. $369-97$; [5] Ч ж а н Л и - ц я н ь, «Математика", 1960, т. 4, № 2, c. $135-59$; [6] Бо p o в k o в A. A., "Изв. АН СССР. Сер. матем.», 1962, т. 26, № 4, c. 605-24.



СМИРНОВА ОБЛА М. С. Никулин. о б ласть т и п а $S,-$ ограниченная односвязная область $G$ с жорданової спрямляемой границей на комплексної шјоскости $\mathbb{C}$ со своӥством: существует такое однолистное гонформное отображение $z=\varphi(w)$ круга $|w|<1$ на область $G$, тто гармонич, функція $\ln \left|\varphi^{\prime}(w)\right|$ при $|w|<1$ представима интегралом Пуассона по своим угловым граничным значениям $\ln \left|\varphi^{\prime}\left(e^{i \theta}\right)\right|$ :



\begin{center}

\includegraphics[max width=\textwidth]{2023_07_09_167802a4647b65445049g-1}

\end{center}



Эти области введены В. И. Смирновым [1] в 1928 при исследовании полноты системы многочленов в С.иирнова классе $E_{2}(G)$. Проблема существования несмирновских областей $\mathrm{c}$ жордановыми спрямляемыми границами была решена М. В. Келдышем и М. А. Лаврентьевым [2], давшими тонкую и сложную конструкцию таких областей и соответствующих отображающих функций $\varphi$ с дополнительным свойством: $\left|\varphi^{\prime}\left(e^{i \theta}\right)\right|=1$ почти для всех $e^{i \theta}$. Основные граничные свойства аналитических функций в круге присущи и функциям, аналитическим в С. о., причем многие из таких свойств справедливы в С. о. и только в них. Примеры С. о. дают жордановы области, границы к-рых суть кривые Јянунова или кусочно ляпуновские кривые с ненулевыми углами.



Лит.: [1] С м и р н о в В. И., «ЖЖ. Ленингр. физ.-матем. об-



\includegraphics[max width=\textwidth, center]{2023_07_09_167802a4647b65445049g-1(1)}

т b e B M. A., "Ann. sci. Ecole norm. supér», 1937, t. 54, p. 138; [3] П р п в а л о в И. И., Граничные свойства аналиттиче ских функций, 2 изд., М.- Л., 1950; [4] Л о в а т е p А., в кн.:

Итоги науки. Математический анализ, т. 10, М., 1973, с. $99-$ Итоги науки. Математический анализ, т. 10, М., 1973, с. 99

$259 ;$ [5] Т у м а р к и н Г. Ц., "Вестн. ЈІенингр. ун-та», 1962, № 13, с. 47-55. Е. П. Долженко.



СНЕДЕКОРА РАСПРЕДЕЛЕНИЕ - см. Фищера F-распределение.



СНОБОЛ - алгоритмический язык, предназначенный для программирования задач обработки символьной информации, т. е. представленной словами в нек-ром алффавите. В литературе по грограммированию такие слова наз. с т р о к а м, или ц е поч к а м и, а оборазующие их буквы - ли т е р а ми. На основе начального варианта С., разработанного в нач. 1960-х гг., было создано несколько версий языка, из к-рых найболее стабильной оказалась версия С.-4. Аналогично ребалу, С. имеет своей теоретич. предпосылкой норяальныле алгорифль А. А. Маркова, в к-рых основной вычислителыной операцией является об̈аружение в слове $A$ вхождения заданного подслова $B$ с последующей заменой этого вхождения на другое слово $C$.



Общая структура программ на С. типична для алгоритмич. языков. Программа имеет вид последовательности операторов (и н с т р у к ц и й) присваивания, сопоставления с образцом, замешения, передачи управления, ввода, вывода и останова. В выражениях могут употребляться как примитивные операции, предикаты и функции, так и функции, определяемые программистом (в т. ч. и рекурсивные). Любая инструкция может быть помечена. Метка в С. может трактоваться так строковая переменная, значением к-рой является инструкция, помеченная этой меткой. Основными типами данных являются строки, целые и действительные числа, имена, образцы. Базисным типом является строка, все остальные данные имеют строковое представление, совпадающее с их способом записи в программе. Однородные данные могут объединяться в массивы и таблицы, шроизвольные данные объединяются в наборы заданной длины и с заданными именами для каждої позиции (п о л я) набора. Имена обозначают переменные, метки, формальные параметры и функции. Строки в С. могут быть любой длины. Переменные не имеют постоянно приписываемого им типа, однако операнды каждой операции или примитивной функции о жи и а ю т данных определенного типа, преобразуя аргументы к ожидаемому типу пибо вылавая сообпение об опибке.



Наиболее характерной операцией С. является сопоставление строки с образцом. О 6 р а 3 е ц - это особос выразкение С., создающее в нек-рой послсдоватсль-





\end{document}